\section{引言}

费曼的部分子分布函数（PDF）\cite{Feynman:1973xc} $f(x)$ 是量子色动力学（QCD）中描述硬包容过程的关键构件。作为累积强子结构非微扰信息的对象，PDF 是格点研究的天然课题。然而，PDF 的直接定义涉及光锥 $z^2=0$ 上双定域算符的矩阵元——这样的间隔在欧氏格点上无法访问。

如何从类空间隔获取信息的想法可追溯到 \mbox{W.~Detmold} 与 D.~Lin 的开创性论文 \cite{Detmold:2005gg}，他们建议对重-轻流的深非弹型欧氏关联函数作格点研究。随后，V.~Braun 与 D.~M\"uller \cite{Braun:2007wv} 建议利用欧氏关联函数提取 $\pi$ 介子分布振幅——另一个在硬排除过程的微扰 QCD 研究中起基本作用的函数 \cite{Radyushkin:1977gp}。以``格点截面''形式使用关联函数的想法近来由 Qiu 及其合作者的论文 \cite{Ma:2014jla,Ma:2017pxb} 所倡导。

上述流关联函数都包含一个连接流顶点的夸克传播子。X.~Ji 的提议 \cite{Ji:2013dva} 避开了这一因子：研究准部分子分布（quasi-PDF）$Q(y,p_3)$，它描述强子动量的空间 $z_3$ 分量 $p_3$ 的分布。虽然与描述强子``plus''动量 $p_+=p_0+p_3$ 分布的费曼 PDF $f(y)$ 不同，但在无限动量极限 $p_3\to \infty$ 下二者与 $f(y)$ 一致。

无论在格点上还是通常的连续时空中，所有类型 PDF 的基本对象都是矩阵元 $M(z,p)$，一般地（即忽略不重要的自旋细节）写作 $\langle  p | \phi (0)  \phi(z)  | p \rangle $。由洛伦兹不变性，它是 Ioffe 时间 $(pz)\equiv -\nu$ \cite{Ioffe:1969kf} 和间隔 $z^2$ 的函数，$M(z,p) \equiv {\cal M} (\nu, -z^2)$。

在（形式的）光锥极限 $z^2=0$ 下，${\cal M} (\nu, 0)$ 对 $\nu$ 的傅里叶变换给出 $f(x)$。在此意义上，Ioffe 时间分布（ITD）${\cal M} (\nu, -z^2)$ 对 $\nu$ 的依赖反映 PDF 的纵向结构。如\mbox{文献 \cite{Radyushkin:2016hsy}} 所示，${\cal M} (\nu, -z^2)$ 的 \mbox{$z^2$ 依赖}决定（QCD 情形为直连接的）横向动量依赖部分子分布（TMD）${\cal F} (x, k_\perp)$ 的 \mbox{$k_\perp$ 依赖}。

由于准 PDF $Q(y,p_3)$ 由 ${\cal M} (z_3 p_3, z_3^2)$ 对 $z_3$ 的傅里叶变换给出，其形状会因通过 ${\cal M} (z_3 p_3, z_3^2)$ 第二个宗量进入的非微扰横向动量效应而畸变。尽管从更一般的视角看 ${\cal F} (x, k_\perp)$ 的 \mbox{$k_\perp$ 依赖}提供强子三维结构的信息，对准 PDF 而言它却是造成 $Q(y,p_3)$ 与 $f(y)$ 之间不希望的差异的麻烦因素；在现有准 PDF 格点计算所能达到的动量下，这一差异非常显著。

为减弱 ITD ${\cal M} (\nu, -z^2)$ 中 $z^2$ 依赖的影响，文献 \cite{Radyushkin:2017cyf} 建议考虑约化 ITD ${\mathfrak M} (\nu, -z^2)$，即 ${\cal M} (\nu, -z^2)$ 与静止系分布 ${\cal M} (0, -z^2)$ 之比。虽然没有第一性原理保证 \mbox{$z^2$ 依赖}的非微扰部分在该比值中消失，但自然可以预期它会大幅减弱。

${\mathfrak M} (\nu, -z^2)$ 只是 $\nu$ 的函数的理想情形对应 TMD ${\cal F} (x, k_\perp)$ 的 $x$ 与 $k_\perp$ 依赖的因子化。事实上，在软区域 $k_\perp^2 \lesssim 1$ GeV$^2$ 内 ${\cal F} (x, k_\perp) = f(x) K(k_\perp)$ 的想法是 TMD 从业者的标准假设（例如见文献 \cite{Anselmino:2013lza}），其中高斯型是 $K(k_\perp)$ 最流行的形式。

文献 \cite{Orginos:2017kos}（另见文献 \cite{Karpie:2017bzm} 的描述）进行了约化 ITD 的探索性格点研究。结果表明 ${\mathfrak M} (\nu, z_3^2)$ 基本上是 $\nu$ 的普适函数，仅对应最小 $z_3$ 值的点与共同曲线有小的偏离。

如文献 \cite{Orginos:2017kos} 所证明，这些偏离可以用微扰演化解释。虽然文献 \cite{Orginos:2017kos} 采用的领头对数近似（LLA）足以分析 $\ln z_3^2$ 依赖，但要指定应赋予所提取的标度相关 PDF $f(x,\mu^2)$ 的标度 $\mu$，就需要超出 LLA。

为此，需要 ITD 单圈修正的完整表达式。最近，文献 \cite{Ji:2017rah,Radyushkin:2017lvu} 报告了此类计算。本文的目标是更详细地讨论演化的 LLA 处理，并把分析扩展到 LLA 之外。我们将证明，单圈修正包含一个大贡献，它相当程度地改变了 LLA 得到的结果。

为使本文自足，我们在\mbox{第 II 节}概述 Ioffe 时间分布与伪 PDF 的基础知识。第 III 节讨论单圈修正的结构。第 IV 节描述文献 \cite{Orginos:2017kos} 格点研究中揭示的演化效应，并把约化 ITD ${\mathfrak M} (\nu, z_3^2)$ 的数据转换为在 $\overline{\rm MS}$ 方案中定义的标准部分子密度 $f(x,\mu^2)$。全文总结与结论见\mbox{第 V 节}。
